% Composition Shift and the Measurement of Bias in a Growing Prediction Market
% Elsevier elsarticle build of reports/final/paper_a_composition.md
%
% NOT COMPILE-VERIFIED. No LaTeX toolchain was available on the machine where
% this was written, so it has been checked structurally (balanced environments,
% escaped specials, every figure present) by tools/lint_tex.py but never run
% through pdflatex. Build it before believing it.
%
% Build:  pdflatex paper_a_composition && pdflatex paper_a_composition
% Needs:  elsarticle.cls, booktabs, amsmath, hyperref (all in TeX Live)

\documentclass[preprint,12pt]{elsarticle}

\usepackage[utf8]{inputenc}
\usepackage[T1]{fontenc}
\usepackage{amsmath}
\usepackage{booktabs}
\usepackage{textcomp}
\usepackage[hidelinks]{hyperref}

\journal{International Journal of Forecasting}

\begin{document}

\begin{frontmatter}

\title{Composition Shift and the Measurement of Bias in a Growing
Prediction Market: Evidence from Kalshi, 2021--2026}

\author{Vladyslav Yurchyna}
\ead{vlad.yurchina@outlook.com}
\address{Independent Researcher}

\begin{abstract}
Prediction markets are now studied intensively enough that their own growth has
become a measurement problem. Using 376.8 million Kalshi fills through June
2026, we show that sports contracts went from 0.009\% of the in-scope sample in
calendar 2024 to 58.8\% after May 2025, and that 41\% of the aggregate change in
the Mincer--Zarnowitz slope at that boundary is composition rather than any
change in how the market prices risk. The natural before/after design across
early 2025 largely measures the arrival of an asset class.

We then ask the persistence question properly. Reconstructing Kalshi's fee
schedule from sixteen dated captures of its own published document, we find the
2025 maker fee was not the exchange-wide regime change the literature assumes: a
per-series surcharge covering 8.0\% to 14.4\% of in-scope markets, revised five
times, flat per contract before it was quadratic in price. That rules out a step
dummy and calls for a difference-in-differences between treated and untreated
series trading in the same months.

The pre-specified estimate is a tight null: $-0.0016$ (0.0064) on 98 treated
series. An event study added after that result was known supports parallel
trends ($\chi^2(5) = 3.24$, $p = 0.66$) and, read as exploratory, indicates a
transitory reduction three to four months after a series is charged. The maker's
return advantage at prices at or above 50c survives every plausible fee rate. We
publish the fee history, the category mapping and the full pipeline.
\end{abstract}

\begin{keyword}
prediction markets \sep favorite--longshot bias \sep composition effects \sep
difference-in-differences \sep transaction fees
\JEL G14 \sep G13 \sep D47 \sep C52
\end{keyword}

\end{frontmatter}

% ---------------------------------------------------------------------------
\section*{Pre-specification}
\addcontentsline{toc}{section}{Pre-specification}

The R2 specification --- the composition decomposition, the static
difference-in-differences, the verdict thresholds --- was committed on
2026-07-28, before any R2 estimate existed. \textbf{The event study in
Section~\ref{sec:eventstudy} was not}; it was added on 2026-08-27, after the
static estimate was known. Section~\ref{sec:whatstatic} states which of its two
halves is a validity check and which is exploratory, and the abstract labels the
exploratory result as such.

% ---------------------------------------------------------------------------
\section{Introduction}

The favorite--longshot bias --- cheap contracts winning less often than their
prices imply, expensive ones more often --- is among the most replicated
findings in the study of betting and prediction markets
\citep{thaler1988,ottaviani2008,snowberg2010}. Kalshi, the first federally
licensed prediction market in the United States, has recently made it newly
measurable at scale: transaction-level data are public, the venue is large, and
its contracts settle unambiguously.

Two features of 2025 make it look like a natural experiment. Kalshi began
charging fees to liquidity providers, taxing the one margin recent work
identifies as exploitable. And in September 2025 the biases themselves were
publicly documented in a widely circulated working paper, whose authors closed
by inviting exactly the follow-up we conduct: whether the patterns persist ``now
that they have been publicly documented.''

The natural design is a before/after comparison of the aggregate bias.
\textbf{Our first result is that this design is contaminated, and we measure by
how much.}

Kalshi did not merely grow across that boundary; it changed shape. Sports
contracts launched in early 2025 --- almost exactly at the boundary any such
study would use --- and by June 2026 they are 58.8\% of the in-scope sample
against 0.009\% in calendar 2024. Because the bias is heterogeneous across
categories, an aggregate comparison mixes a change in pricing behaviour with a
change in what is being priced. Decomposed, 41\% of the aggregate movement at
the fee boundary is the composition term.

This is not a hypothetical concern. Kalshi is being analysed now, on samples
that straddle the same boundary, and a naive aggregate comparison will report a
large apparent effect.

Our remaining contributions follow from taking the confound seriously.

\paragraph{The fee treatment is not what it is assumed to be.} We reconstruct
Kalshi's fee schedule from sixteen dated captures of the exchange's own
published document, 2021-07 through 2026-02, plus per-series fee metadata for
all 12,231 series. The maker fee was never exchange-wide: it is a surcharge on
an enumerated list of series, 29 at introduction and 111 by September 2025,
covering 8.0\% to 14.4\% of in-scope markets and 61\% to 66\% of volume. Its
first form was flat per contract, becoming quadratic in price only on
2025-07-08. Two further carve-outs --- a 0.14 taker rate before August 2021, and
a half rate for index markets covering 18.3\% of the pre-2025 universe --- are
not modelled in the existing literature at all.

\paragraph{Properly identified, the fee produced no persistent change in the
bias.} Because treatment is per-series, staggered and revised, treated and
untreated series trade in the same months and a difference-in-differences is
available. The pre-specified estimate is $-0.0016$ (s.e.\ 0.0064) against
never-treated series. Treatment here was selected on market-making activity,
which is adjacent to the outcome, so the identifying assumption needs evidence
rather than assertion; an event study added after that estimate was known finds
pre-treatment coefficients jointly zero ($p = 0.66$), which supports it.

That same event study also finds post-treatment coefficients jointly non-zero
($p = 0.021$), suggesting a transitory reduction peaking three to four months
after a series is charged. \textbf{We report this as exploratory}, for the
reason Section~\ref{sec:whatstatic} sets out: it is a hypothesis test run after
a pre-specified test returned zero, which is the shape specification search
produces even when nobody went looking. The naive step-dummy coefficient
disagrees with both, for a third reason again: it is dominated by one category
that has shrunk by a factor of fifty while retaining nearly a third of the
composition-held-fixed weight.

\paragraph{The exploitable margin survives.} The maker's return advantage at
prices at or above 50c is $+2.40\%$ gross in the post-boundary window and never
crosses zero across the plausible band of maker fee rates.

Section~\ref{sec:data} describes the data. Section~\ref{sec:validation}
validates the pipeline against published results.
Section~\ref{sec:composition} is the composition result. Section~\ref{sec:fees}
is the fee schedule. Section~\ref{sec:persistence} is the persistence estimate.
Sections~\ref{sec:robustness} and~\ref{sec:related} cover robustness and related
work.

% ---------------------------------------------------------------------------
\section{Data}\label{sec:data}

All data come from Kalshi's public \texttt{trade-api/v2} endpoints --- no
account, no authentication, no order placement. Collection runs in two passes: a
universe-wide pass for market metadata, the boundary trades a daily-lookback
price panel requires, and closing quotes for the spread filter; then a full-tape
pass restricted to contracts surviving the filters, because maker/taker
quantities need every fill's price, size and taker side.

The completed tape is \textbf{376,760,957 fills across 62 monthly partitions}.
The in-scope set for the post-boundary window --- markets clearing \$1,000
volume, 24 hours open, and a closing bid--ask spread at or under 20c --- is
\textbf{391,427 markets}, all of them taped.

Three construction choices are stated because each silently changes the sample,
and because the second is a correction to the natural reading of the prior
literature.

\paragraph{Volume is denominated in contracts, not dollars.} The original screens
on ``We focus only on contracts that have reached a total trading volume upon
market closure of at least \$1,000'', which reads as dollar notional, but
Kalshi's \texttt{volume} field counts contracts and the API exposes no notional
field. At
the contract reading our independently collected universe lands within 0.1\% of
the published event count; the dollar reading roughly halves the sample. Since a
sample rule that differentially removes cheap strikes is itself a result in a
paper about the favorite--longshot bias, both are computed and the contract
reading is primary.

\paragraph{Prices are carried forward on no-trade lookback days.} The original
collects ``the final traded price as the market closed and also, where
available, previous prices from 24-hour intervals up to 10 days before markets
closed'' --- and ``where available'' reads as skipping empty days. The published
counts settle it the other way: 156,986 prices over 46,282 contracts
is 3.39 per contract, which skipping cannot reach. On our tape, skipping yields
2.50 and carrying forward yields 3.50. Carrying forward is primary.

\paragraph{One construction on both sides of every boundary.} The estimands
below measure a \emph{change} in slope within our own data against our own
reference bias. A construction that switched at a boundary is the one thing that
would break identification, so the rule is fixed before any estimate is
computed.

Categories come from Kalshi's own taxonomy, mapped through a versioned table in
the repository.

% ---------------------------------------------------------------------------
\section{Validation}\label{sec:validation}

Before using this pipeline to make a claim about a boundary, we check that it
reproduces what is already known about the pre-boundary window.
\citet{burgi2026}, hereafter BDW, provide the reference points.

\begin{table}[htbp]
\centering
\caption{Reproduction of BDW's reference quantities.}
\begin{tabular}{lrr}
\toprule
Quantity & Ours & BDW \\
\midrule
Maker return, price $\geq$50c & $+2.09\%$ & $+2.6\%$ \\
Taker return                  & $-28.40\%$ & $-31.46\%$ \\
Maker share, 1--10c band      & 43.1\% & 43.5\% \\
Maker share, 90--99c band     & 57.0\% & 56.5\% \\
$\psi$, 2021                  & 0.0409 & 0.041 \\
$\psi$, 2025 (Jan--Apr)       & 0.0171 & 0.021 \\
\bottomrule
\end{tabular}
\end{table}

All five of their by-year slopes are recovered with matching sign and
significance pattern; magnitudes run slightly below theirs, within roughly one
standard error. The bias itself reproduces directly: in the 1--10c band
contracts win 2.43\% of the time at a mean price of 3.22c, and in the 91--99c
band 97.86\% at 97.08c.

Our sample is smaller than theirs --- 33,222 contracts against 46,282 --- and
the shortfall is almost entirely the spread filter: 3,257 markets (8.2\% of
those clearing volume and duration) have no bid--ask history that Kalshi will
serve, so the loss is structural rather than a collection gap. The full
reconciliation, the construction branches and the divergences are the subject of
the companion replication and are not repeated here.

Two checks worth stating. Kalshi's taker-side fields are populated on
\textbf{100.0\%} of all 376.8 million fills in every era. And the two-way
(event, contract) cluster-robust variance reduces to one-way event clustering
when contracts nest in events \citep{cameron2011}; estimated both ways on the
real panel the standard errors agree to fourteen decimal places, so one-way
event clustering is used throughout without hedging.

% ---------------------------------------------------------------------------
\section{Composition}\label{sec:composition}

\subsection{The venue changed shape, not just size}

Kalshi launched sports contracts in early 2025. The timing is unfortunate for
anyone studying the 2025 fee change, because it is essentially the same date.

\begin{table}[htbp]
\centering
\caption{Category weights, calendar 2024 against the post-boundary window.}
\begin{tabular}{lrrr}
\toprule
Category & Calendar-2024 weight & Post-boundary weight & Change \\
\midrule
\textbf{Sports}     & \textbf{0.00009} & \textbf{0.5875} & \textbf{$+0.5874$} \\
Climate \& Weather  & 0.4744 & 0.0995 & $-0.3749$ \\
Financials          & 0.3047 & 0.0059 & $-0.2989$ \\
Mentions            & 0.0000 & 0.0646 & $+0.0646$ \\
Economics           & 0.0690 & 0.0086 & $-0.0604$ \\
\bottomrule
\end{tabular}
\end{table}

Four categories present after the boundary --- Exotics, Mentions, Social,
Transportation --- did not exist in calendar 2024 at all.

This matters because the slope is category-heterogeneous. Estimated separately,
$\psi$ ranges from 0.0107 in Entertainment to 0.0743 in Sports and 0.1206 in
Transportation. Re-weighting a heterogeneous quantity is not a neutral act.

\subsection{Decomposing the aggregate change}\label{sec:decomp}

Write the aggregate change in slope at a boundary as
\begin{equation}
\Delta\psi_{\mathrm{agg}}
  = \underbrace{\sum_c \bar{w}_c \, \Delta\psi_c}_{\text{within}}
  + \underbrace{\sum_c \Delta w_c \, \bar{\psi}_c}_{\text{between}},
\end{equation}
where the first term holds composition fixed at the calendar-2024 mix and varies
only within-category behaviour, and the second holds behaviour fixed and varies
only weights. At the fee boundary:

\begin{table}[htbp]
\centering
\caption{Within/between decomposition at the fee boundary.}
\begin{tabular}{lrr}
\toprule
Term & Value & Share \\
\midrule
Within (behaviour)               & $+0.0399$ & 59\% \\
\textbf{Between (composition)}   & \textbf{$+0.0278$} & \textbf{41\%} \\
Aggregate                        & $+0.0677$ & 100\% \\
\bottomrule
\end{tabular}
\end{table}

The between term is dominated by one category. Sports alone contributes
$+0.0437$ --- more than the entire between component, with the shrinking
categories contributing negatively against it.

\subsection{What this implies}

An aggregate before/after comparison across early 2025 on Kalshi will find a
large apparent change in bias, of which roughly two fifths is the arrival of
sports contracts. The direction happens to be positive here, which is worth
noting: a naive reading of the aggregate would report that the bias
\emph{strengthened} after a fee designed to tax it --- a conclusion with no
plausible mechanism behind it, arrived at by not asking what the sample was made
of.

\subsection{What is and is not new here}

The general principle is old. That an aggregate can move in the opposite
direction to every subgroup composing it is \citet{simpson1951}; separating a
between-group weight change from a within-group behavioural change is the
\citet{oaxaca1973} and \citet{blinder1973} decomposition, and the shift-share
logic that underlies \citet{bartik1991} instruments. We claim no novelty for the
arithmetic in Section~\ref{sec:decomp} and we use the standard form
deliberately, so that the object being reported is familiar.

What is specific to this setting is worth stating plainly, because the obvious
referee response to Section~\ref{sec:decomp} is that it restates a textbook
fact.

\paragraph{The magnitude is not a textbook magnitude.} Shift-share
decompositions were developed for panels whose composition drifts over decades.
A single category going from 0.009\% to 58.8\% of the sample in fourteen months
is not drift; it is the venue becoming a different venue. The resulting
distortion is large enough to invert the economic reading rather than merely
attenuate it: the naive aggregate says a fee levied on liquidity providers made
the market \emph{less} efficiently priced, which is not a mechanism anyone would
defend if it were stated aloud.

\paragraph{The original's own robustness check does not reach this.} BDW
anticipate the category question and report that ``Our results below are not
sensitive to cutting the data off in December 2024 or to excluding the category
containing sports bets.'' That check is correct on their sample and does not
cover ours, for a reason that is arithmetic rather than methodological: sports
is 	extbf{4.3\% of the observations inside their window} and 	extbf{58.8\%
after the boundary}, a factor of fourteen. Dropping a category worth one
observation in twenty-three barely moves an estimate, which is what they found.
The problem we identify is created by the share rising to a majority
\emph{after} their sample ends, so no test run inside their window could have
detected it. Their check also asks a different question --- whether a level
estimate survives deleting a category --- where ours asks how much of a
\emph{change across a boundary} is reweighting. Deleting sports and holding the
mix fixed are not the same operation.

\paragraph{The trap is live, not hypothetical.} Kalshi is being analysed now,
and the samples in circulation straddle exactly this boundary ---
\citet{becker2026} covers June 2021 to November 2025, which contains the sports
launch in the middle. We are not warning about a possible future error.

\paragraph{And it is a property of young markets specifically.} A mature
exchange adds listings within existing asset classes; a new one adds asset
classes. Prediction markets are currently doing the latter --- Kalshi added
sports in 2025, and Polymarket's own composition has moved comparably --- so any
study of this literature's growth period inherits the problem rather than
encountering it occasionally.

The contribution is therefore the measurement and its consequence in a setting
where several groups are actively estimating the affected quantity, not the
decomposition identity.

The remedy we adopt is to fix the mix at the last full pre-sports year and to
report the two terms separately, with every behavioural claim referring to the
within term only. Sports is then reported as its own stratum rather than being
allowed to dominate an aggregate it did not exist to generate. Any equivalent
strategy would do; the point is that some strategy is required, and that its
absence is not visible in an aggregate coefficient.

% ---------------------------------------------------------------------------
\section{What the fee schedule actually says}\label{sec:fees}

The design in Section~\ref{sec:persistence} depends on the fee treatment being
what we say it is, so we sourced it rather than assuming it. Sixteen dated
captures of Kalshi's own published fee document (2021-07 to 2026-02) are
archived in the repository with the parse that produces the machine-readable
schedule, alongside per-series fee metadata for all 12,231 series obtained from
Kalshi's API.

\paragraph{The maker fee was never exchange-wide.} Resting orders are free, in
the document's own words, ``unless they are included in our `Maker Fees'
section'' --- an enumerated list of 29 series at introduction, then 39, 48, 101
and 111 by 2025-09-17. On our in-scope universe that is 8.0\% of markets and
61.0\% of volume in the first-list era, rising to 14.4\% and 66.2\%. A minority
of markets and a majority of dollars --- which is why the treatment is
economically real and statistically awkward at the same time.

\paragraph{It was flat before it was quadratic.} The first form is
$\mathrm{round\_up}(0.0025 \cdot C)$ per contract, with no price dependence
whatsoever, becoming $0.0175 \cdot C \cdot P(1-P)$ only on 2025-07-08.

\paragraph{It began on 2025-05-13}, by the document's own revision stamp.

\paragraph{Two carve-outs are absent from existing models.} The taker rate was
0.14 in the July 2021 capture, dropping to 0.07 on 2021-08-01. And S\&P 500 and
Nasdaq-100 markets pay half rate --- 0.035 --- in every capture from September
2022 onward, covering \textbf{18.3\%} of the pre-2025 in-scope universe.

One gap is open and carried as bounds rather than closed by assumption: from the
2025-10-01 revision the document stopped enumerating series, so membership after
that date is propagated through the fee-sensitivity analysis as evidence bounds
(103 series primary, 138 upper).

The methodological consequence is larger than the individual corrections. A step
dummy at a single exchange-wide date is the wrong estimator for a treatment that
touched a minority of markets, was revised five times, and whose intensity ran
5.5\% $\rightarrow$ 9.4\% $\rightarrow$ 17.6\% $\rightarrow$ 14.5\%
$\rightarrow$ 32.0\% $\rightarrow$ 33.5\% before decaying. But the same facts
supply the remedy: because treatment is per-series and staggered, treated and
untreated series trade side by side in the same months.

% ---------------------------------------------------------------------------
\section{Persistence, properly identified}\label{sec:persistence}

\subsection{The difference-in-differences}

We estimate
\begin{align}
(Y-P) ={}& \alpha + \psi P + \alpha_E \mathrm{Ever} + \delta_E (\mathrm{Ever} \cdot P) \nonumber \\
        &+ \sum_m \left[ \alpha_m D_m + \delta_m (D_m \cdot P) \right]
         + \alpha_D \mathrm{Now} + \delta_D (\mathrm{Now} \cdot P) + \varepsilon,
\end{align}
where \textrm{Now} indicates that the market's series carried a maker fee on the
day it closed, \textrm{Ever} marks every observation of a series treated at some
point, and $D_m$ are calendar-month dummies entering the slope as well as the
level, because the estimand is a slope. $\delta_D$ is identified from variation
\emph{within a month, between} series that carried the fee and series that did
not. Treatment is read from the sourced schedule through the same lookup used
for net returns, so the design and the fee accounting cannot disagree about who
was treated.

\begin{table}[htbp]
\centering
\caption{The maker-fee difference-in-differences.}
\begin{tabular}{lrrr}
\toprule
Specification & $\delta_{\mathrm{did}}$ & s.e. & 95\% CI \\
\midrule
Two-way fixed effects & $+0.0113$ & 0.0129 & $[-0.0139, +0.0365]$ \\
\textbf{Clean controls} & \textbf{$-0.0016$} & \textbf{0.0064} & $[-0.0141, +0.0109]$ \\
\bottomrule
\end{tabular}
\end{table}

Estimated on 98 treated series, 141,708 treated observations against 1,188,737
controls, 119,646 event clusters across 60 months, with no cell thin enough to
require a wild cluster bootstrap.

Because two-way fixed effects under staggered adoption permits already-treated
units to serve as controls for later-treated ones, with the attendant negative
weighting \citep{goodmanbacon2021}, both fits are reported always; the
clean-controls variant restricts comparisons to never-treated series --- the
correction \citet{callaway2021} and \citet{sun2021} formalise --- and is the one
reported on disagreement.

\textbf{Read on its own, this says the maker fee did not change the average
slope. Section~\ref{sec:eventstudy} shows that reading is incomplete.}

\subsection{The identifying assumption, and why selection sharpens it}

$\delta_{\mathrm{did}}$ is unbiased only under \textbf{parallel trends}: absent
the fee, treated and untreated series would have moved together in slope. This
is an assumption, not a result, and it deserves more scrutiny here than it
usually gets, because assignment to treatment was plainly not random.

Kalshi charged an enumerated list of series. Section~\ref{sec:fees} shows which:
8.0\% to 14.4\% of in-scope markets carrying 61\% to 66\% of volume --- that is,
the series where market-making is most active and most profitable. Maker
profitability is not incidental to what we are measuring. It is economically
linked to the very slope the estimand concerns:
Section~\ref{sec:margin} reports the maker's return advantage as a headline
quantity, and an exchange choosing to tax the series where that advantage is
largest is selecting on something adjacent to the outcome.

The \textrm{Ever} terms absorb \textbf{permanent} differences in level and slope
between the charged and uncharged groups, and those differences are certainly
present. They do not absorb a \textbf{differential trend}. Selection of this
kind therefore raises the bar that parallel trends must clear rather than
lowering it, and asserting the assumption without evidence would be the weakest
point in the paper.

\subsection{Event study: pre-trends, and a dynamic effect the average
hides}\label{sec:eventstudy}

We re-estimate with leads and lags measured in months from \textbf{each series'
own first treatment month}, never-treated series anchoring the calendar path,
$k=-1$ omitted as the reference and endpoints binned. The sample is the same
1,330,445 observations and 119,646 clusters; 98 treated series fall into 11
adoption cohorts.

\begin{table}[htbp]
\centering
\caption{Event-study coefficients around each series' own adoption month.}
\begin{tabular}{rrrr}
\toprule
Event time $k$ & $\delta_k$ & s.e. & 95\% CI \\
\midrule
$-6$ & $-0.0428$ & 0.0285 & $[-0.0988, +0.0131]$ \\
$-5$ & $-0.0187$ & 0.0624 & $[-0.1410, +0.1037]$ \\
$-4$ & $-0.0975$ & 0.0875 & $[-0.2690, +0.0740]$ \\
$-3$ & $+0.0016$ & 0.0633 & $[-0.1225, +0.1257]$ \\
$-2$ & $-0.0235$ & 0.0410 & $[-0.1039, +0.0569]$ \\
$\mathbf{-1}$ & --- & --- & \emph{reference} \\
0    & $-0.0092$ & 0.0258 & $[-0.0598, +0.0414]$ \\
$+1$ & $+0.0179$ & 0.0278 & $[-0.0365, +0.0724]$ \\
$+2$ & $-0.0120$ & 0.0291 & $[-0.0691, +0.0451]$ \\
$\mathbf{+3}$ & $\mathbf{-0.0771}$ & 0.0342 & $\mathbf{[-0.1441, -0.0102]}$ \\
$\mathbf{+4}$ & $\mathbf{-0.0826}$ & 0.0342 & $\mathbf{[-0.1496, -0.0157]}$ \\
$+5$ & $-0.0462$ & 0.0369 & $[-0.1185, +0.0261]$ \\
$+6$ & $-0.0210$ & 0.0257 & $[-0.0713, +0.0293]$ \\
$+7$ & $-0.0552$ & 0.0346 & $[-0.1230, +0.0127]$ \\
$+8$ & $-0.0242$ & 0.0257 & $[-0.0746, +0.0261]$ \\
$+9$ & $-0.0037$ & 0.0244 & $[-0.0515, +0.0441]$ \\
\bottomrule
\end{tabular}
\end{table}

\paragraph{Parallel trends is supported.} Every pre-treatment interval covers
zero, and the joint Wald test that all five are zero gives $\chi^2(5) = 3.24$,
$p = 0.66$. Slopes were not already diverging before Kalshi charged anybody.
Given the selection argument above, this is the evidence the design needed
rather than a formality.

\paragraph{The post-treatment coefficients are jointly non-zero --- an
exploratory finding, labelled as such here and everywhere it is repeated.} The
joint test that all ten are zero gives $\chi^2(10) = 21.01$, $p = 0.021$,
rejecting at 5\%. We report the joint test rather than the two bolded rows,
because ten coefficients tested individually at 5\% would produce at least one
``significant'' month roughly forty percent of the time under the null; the
joint statistic is what establishes there is something to explain.

The exploratory label is not modesty. This test was not pre-specified and was
run after the pre-specified one returned zero --- see the status note at the end
of this section, which readers should take as attached to every sentence in the
next two paragraphs.

The shape is coherent: nothing at impact, a \textbf{negative} deviation peaking
at three and four months after a series is charged --- the favorite--longshot
slope \emph{falling} in treated series relative to controls --- and decay back
toward zero by nine months. So the fee did move the bias, transitorily and with
a lag, in the direction of less bias.

We flag two things we cannot settle. The delayed onset is not obviously
explicable by a fee that takes effect immediately, and we do not have a
mechanism we would defend; candidate explanations include contracts written
before treatment settling after it, and adjustment in maker behaviour rather
than in posted prices. And the point estimates at $k=+3$ and $+4$ are large
relative to the baseline slope of 0.0254 --- they imply a temporary sign
reversal --- with correspondingly wide intervals. We report the shape and the
joint test, and we do not build a magnitude claim on two coefficients.

\paragraph{When this analysis was added, and why the two halves differ in
status.} The static difference-in-differences was pre-specified in the analysis
plan committed on 2026-07-28. The event study was not: it was added on
2026-08-27, after the static estimate was known, in response to a critique of
the then-unstated parallel-trends assumption. The two halves of it do not have
the same epistemic standing.

The pre-period test is a \textbf{validity check on an identifying assumption}.
It was worth running whatever it showed, it can only undermine the design rather
than manufacture a finding, and we would have reported a failure as readily as
this pass. Adding such a check late costs nothing but a date.

The post-period joint test is a \textbf{hypothesis test conducted after a
pre-specified test returned zero}, and it returned $p = 0.021$. We label it
exploratory on that basis. We say so plainly because a transitory effect
surfacing in a second analysis, after the first found nothing, is precisely the
pattern that innocent specification search produces --- and a paper arguing that
aggregates hide structure has no standing to be casual about the order in which
its own analyses were run.

Accordingly the pre-specified null remains the headline result for the fee
question, and the hump is reported as a labelled exploratory finding that
warrants a pre-specified test on future data rather than as an established
effect.

\subsection{What the static estimate does and does not say}\label{sec:whatstatic}

The static $\delta_{\mathrm{did}}$ of $-0.0016$ (0.0064) is not wrong; it is an
average, and its tight interval is a statement about that average rather than
about every month within it. Averaging a delayed, transitory, decaying effect
over ten post-treatment months is exactly how it becomes indistinguishable from
zero.

This is worth naming, because it is the same failure mode as
Section~\ref{sec:composition} in a different guise. There, an aggregate over
categories hid the fact that its movement was composition. Here, an aggregate
over event time hides the fact that a null is made of a hump. In both cases the
aggregate is not false, and in both cases reporting only the aggregate would
have been.

The honest summary of the fee question, with each part carrying its status:
\textbf{pre-specified, no persistent change in the bias; exploratory, a
transitory reduction around three to four months after treatment; and separately
established, no evidence that the maker's exploitable margin was eliminated} ---
the last of which Section~\ref{sec:margin} shows directly and which was
pre-specified as the fee-sensitivity ribbon.

\paragraph{A limitation we state rather than bury.} With 11 adoption cohorts,
the event-study coefficients can still be contaminated across cohorts in the way
\citet{sun2021} describe, since a treated cohort's post-period overlaps
another's pre-period. Never-treated series carry most of the identifying weight
here, which limits the problem but does not eliminate it. A fully
interaction-weighted or Callaway--Sant'Anna estimator would settle it, and we
regard that as the natural next step rather than something this draft has done.

\subsection{The naive coefficient, and why it disagrees}

The composition-weighted step coefficient at the fee boundary is $+0.0399$ with
a 95\% interval of $[+0.0071, +0.0726]$ --- nominally significant, positive, and
therefore implying the bias strengthened under a fee designed to tax it.

Two features explain it and neither is behavioural. First, as
Section~\ref{sec:fees} establishes, the step is applied to a sample most of
whose observations were never treated, against a treatment whose intensity
quadrupled and then halved inside the window. Second, its within component is
essentially one category: Financials contributes $+0.0437$ of the $+0.0399$, on
a within-category slope change of $+0.1435$, while holding 30.5\% of the
composition-held-fixed weight against 0.59\% of the current sample. Holding the
mix fixed is correct --- it is the whole point --- but the consequence is that
the headline rests on a cell that has shrunk fiftyfold, and we decline to build
a persistence claim on it.

\subsection{The exploitable margin survives}\label{sec:margin}

The maker's advantage at prices at or above 50c, in the post-boundary window,
computed in three fee layers on 202,725,609 maker and 167,925,545 taker sides
with no observation lost to gaps in the fee history:

\begin{table}[htbp]
\centering
\caption{The maker's advantage at prices at or above 50c, by fee layer.}
\begin{tabular}{lr}
\toprule
Layer & Margin \\
\midrule
Gross & \textbf{$+2.40\%$} \\
Net of own-era fees & $+4.39\%$ \\
Pre-2025 schedule held constant & $+4.54\%$ \\
\bottomrule
\end{tabular}
\end{table}

The net layer exceeding the gross layer is the expected direction rather than an
anomaly: under fees the taker pays and the maker mostly does not, so the spread
between them widens. Swept across an eleven-point grid spanning the plausible
maker fee band, the margin declines monotonically from $+0.0454$ to $+0.0431$
and \textbf{never crosses zero}. There is no break-even rate inside the band.

The margin documented before the boundary is therefore still present after it,
and no plausible maker fee erases it.

\subsection{The publication boundary}

The publication coefficient is $-0.0311$ with an interval of $[-0.0637,
+0.0014]$. It contains zero, and it also contains the negative of our reference
bias. On this sample, persistence and disappearance cannot be distinguished. We
report that rather than choosing.

% ---------------------------------------------------------------------------
\section{Robustness}\label{sec:robustness}

\paragraph{Horizon.} The panel pools eleven daily lookback horizons. Estimated
separately, the naive fee coefficient swings from $+0.098$ at one day to
$-0.008$ at nine; the one-observation-per-contract closing-price specification
gives $+0.0117$ against the pooled $+0.0399$. The difference-in-differences is
not subject to this sensitivity, but the failure of the naive coefficient to
survive the horizon check is a further reason not to read it behaviourally.

\paragraph{Control venue.} Polymarket's monthly slope over 2025-05 to 2025-12
provides a secular-trend check --- not a difference-in-differences, and we do
not use parallel-trends language. Three caveats are stated here rather than left
for a referee. Polymarket's tail bias runs in the opposite direction to Kalshi's
\citep{qin2026}, so the comparison controls for market-wide drift in efficiency
and not for level or sign. The Polymarket data come from a public third-party
archive, which the author worked with in separate research on Polymarket
forecasting --- the project this paper's question originated in. That work
measures forecasters' skill; this one measures the market's own calibration
path, so the archive is put to a different use here, and its full provenance is
documented in the repository.
And cross-venue arbitrage could propagate the publication event to Polymarket as
well, which would strengthen a secular-trend reading rather than undermine the
design.

The control window closes at 2025-12-31 for a structural reason: Polymarket
began its own staggered fee reform in January 2026, so a longer archive would
carry that venue's own treatment inside ours. The fee estimate is unaffected,
being identified within Kalshi. \textbf{The publication estimate is controlled
only through 2025-12-31}, and the remaining six months are reported as an
uncontrolled extension of the horizon rather than averaged into a single figure.

% ---------------------------------------------------------------------------
\section{Related work}\label{sec:related}

\paragraph{Kalshi and the favorite--longshot bias.} \citet{burgi2026} is the
immediate antecedent and remains a working paper as of this draft. Our companion
paper is a full replication of it. The favorite--longshot bias itself is among
the most replicated regularities in betting and prediction markets --- surveyed
by \citet{ottaviani2008}, documented in parimutuel betting by
\citet{thaler1988} and estimated structurally by \citet{snowberg2010} --- and we
take its existence as established rather than as something this paper
demonstrates. On prediction markets as forecasting instruments generally we
follow \citet{wolfers2004}.

\citet{whelan2023} introduces fees into a prediction-market model and shows that
fees levied on winnings generate a form of favorite--longshot bias in post-fee
loss rates. It is theoretical and uses no Kalshi data. Section~\ref{sec:fees}
speaks directly to its assumptions: the fee it models as uniform is neither
uniform nor, in its first eight months, price-dependent.

\citet{becker2026} analyses 72.1 million Kalshi trades from June 2021 to
November 2025, confirms the longshot bias by price level, decomposes returns by
market role --- takers $-1.12\%$ and makers $+1.12\%$ mean excess return per
trade --- and documents a YES/NO asymmetry in taker order flow. Its role
decomposition agrees with ours in direction; magnitudes are not comparable,
since the estimand is excess return per trade rather than return on a contract
position. It does not address the fee regime, dated boundaries, or composition,
and its sample ends seven months before ours. We note that its window straddles
the sports launch, which is the configuration Section~\ref{sec:composition} is
about.

\paragraph{Anomalies after publication.} The publication boundary asks a
question with an established finance literature behind it. \citet{mclean2016}
estimate that cross-sectional return predictors decay by roughly 58\% out of
sample after publication, of which about a third is attributable to publication
itself rather than to in-sample overfitting, and interpret the residual decay as
informed trading against a documented signal. The prediction market case is
cleaner in one respect and harder in another: contracts settle unambiguously, so
there is no benchmark model to misspecify, but a single venue with two candidate
treatments four months apart supports far less separation than a cross-section
of thousands of predictors. BDW's own closing invitation is this literature's
question asked of their result, and we report that our sample cannot answer it.

\paragraph{Staggered difference-in-differences.} Our treatment is adopted at
different dates by different series, which is the setting in which conventional
two-way fixed effects is now known to misbehave. \citet{goodmanbacon2021} shows
that the TWFE estimand is a weighted average of all available two-group
comparisons, including ones that use already-treated units as controls for
later-treated ones, and that those comparisons can receive negative weight when
effects are heterogeneous over time. \citet{callaway2021}, \citet{sun2021} and
\citet{dechaisemartin2020} each propose estimators that restrict comparisons to
clean control groups and aggregate the resulting group-time effects explicitly.

Our clean-controls specification is the simplest member of that family: it
restricts every comparison to treated-versus-never-treated, which is precisely
the correction those papers identify as necessary. We report it alongside the
TWFE fit rather than instead of it, and treat their agreement as the evidence
that weighting is not driving the result. Where they disagree, the
clean-controls estimate is the one reported. The event study in
Section~\ref{sec:eventstudy} is specified in the same spirit --- event time is
measured relative to each series' own adoption date, and never-treated series
anchor the calendar path.

\paragraph{Composition and decomposition.} The within/between decomposition in
Section~\ref{sec:decomp} is standard: \citet{simpson1951} for the aggregation
reversal, \citet{oaxaca1973} and \citet{blinder1973} for the two-term form, and
the shift-share tradition behind \citet{bartik1991} for holding a mix fixed
while letting behaviour vary.

% ---------------------------------------------------------------------------
\section{Conclusion}

A venue that grows by changing what it lists cannot be studied with an aggregate
before/after comparison. On Kalshi, 41\% of the apparent change in bias at the
2025 fee boundary is the arrival of sports contracts, which went from 0.009\% to
58.8\% of the sample in fourteen months. The remedy is cheap --- fix the mix,
decompose, report both terms --- and its absence is invisible in the coefficient
it distorts.

On the question that motivated the exercise: the maker fee, correctly measured
as the per-series and repeatedly revised surcharge it actually was, produced
\textbf{no persistent change in the bias}. That is the pre-specified result and
it stands as the answer.

An event study added afterwards, in response to a referee's question about the
identifying assumption, supports that assumption --- pre-treatment slopes were
not diverging, which the selection into treatment made worth checking rather
than assuming --- and additionally suggests a transitory reduction peaking three
to four months after a series is charged. We label that second finding
exploratory and do not treat it as established, because it emerged from an
analysis run after the pre-specified one returned zero. If it is real, it should
survive a pre-specified test on data this paper does not yet have, and that is
the test we would want run before anyone believes it.

The maker's advantage at prices at or above 50c persists after the boundary and
survives every plausible fee rate, so whatever the fee did, it did not remove
the edge. Whether the public documentation of these patterns changed them cannot
be settled on this sample, and we do not claim it can.

The fee history, the category mapping, the archived source documents and the
complete pipeline are public.

% ---------------------------------------------------------------------------
\section*{Declaration of interest}

The author declares no competing financial interest. The author does not trade
on Kalshi, holds no positions on the venue studied, has no financial or personal
relationship with the exchange or with the authors of the work replicated, and
received no funding for this research.

\section*{Data and code availability}

The repository contains the collection pipeline, the analysis code, sixteen
archived captures of Kalshi's fee schedule with the parser that turns them into
a machine-readable step function, the category mapping table, and the committed
analysis plan with its dated amendments. It is MIT-licensed and contains no
execution code: every endpoint used is public and unauthenticated.

% Elsevier requires this as its own section, immediately before the reference
% list. It covers the writing process only -- it is deliberately not a
% methodology statement, and nothing in this paper's analysis treats AI as an
% object of study.
\section*{Declaration of Generative AI and AI-assisted technologies in the
writing process}

During the preparation of this work the author used Claude (Anthropic) in order
to draft and revise manuscript prose and to generate the data-collection and
analysis code in the accompanying repository. After using this tool the author
reviewed and edited the content as needed and takes full responsibility for the
content of the publication.

% ---------------------------------------------------------------------------
\begin{thebibliography}{99}

\bibitem[Bartik(1991)]{bartik1991}
Bartik, T.J., 1991.
\emph{Who Benefits from State and Local Economic Development Policies?}
W.E. Upjohn Institute for Employment Research, Kalamazoo, MI.

\bibitem[Becker(2026)]{becker2026}
Becker, J., 2026.
The Microstructure of Wealth Transfer in Prediction Markets.
Self-published 18 January 2026; formalized version, SSRN Working Paper 7217640,
1 August 2026.

\bibitem[Blinder(1973)]{blinder1973}
Blinder, A.S., 1973.
Wage discrimination: reduced form and structural estimates.
\emph{Journal of Human Resources} 8 (4), 436--455.

\bibitem[B\"urgi et al.(2026)]{burgi2026}
B\"urgi, C., Deng, W., Whelan, K., 2026.
\emph{Makers and Takers: The Economics of the Kalshi Prediction Market.}
Working paper, January 2026 version. CEPR Discussion Paper 20631; CESifo Working
Paper 12122; UCD Working Paper WP2025\_19; GWU Working Paper 2026-001; MPRA
Paper 126350. (The GWU series cover sheet, dated February 2026, reads ``Makers
\emph{or} Takers''; the paper's own title page inside that document reads
``Makers \emph{and} Takers'' and its 44 pages are identical to the January
version. We cite the authors' title.)

\bibitem[Callaway and Sant'Anna(2021)]{callaway2021}
Callaway, B., Sant'Anna, P.H.C., 2021.
Difference-in-differences with multiple time periods.
\emph{Journal of Econometrics} 225 (2), 200--230.

\bibitem[Cameron et al.(2011)]{cameron2011}
Cameron, A.C., Gelbach, J.B., Miller, D.L., 2011.
Robust inference with multiway clustering.
\emph{Journal of Business \& Economic Statistics} 29 (2), 238--249.

\bibitem[de Chaisemartin and d'Haultf{\oe}uille(2020)]{dechaisemartin2020}
de Chaisemartin, C., d'Haultf{\oe}uille, X., 2020.
Two-way fixed effects estimators with heterogeneous treatment effects.
\emph{American Economic Review} 110 (9), 2964--2996.

\bibitem[Goodman-Bacon(2021)]{goodmanbacon2021}
Goodman-Bacon, A., 2021.
Difference-in-differences with variation in treatment timing.
\emph{Journal of Econometrics} 225 (2), 254--277.

\bibitem[McLean and Pontiff(2016)]{mclean2016}
McLean, R.D., Pontiff, J., 2016.
Does academic research destroy stock return predictability?
\emph{Journal of Finance} 71 (1), 5--32.

\bibitem[Oaxaca(1973)]{oaxaca1973}
Oaxaca, R., 1973.
Male--female wage differentials in urban labor markets.
\emph{International Economic Review} 14 (3), 693--709.

\bibitem[Ottaviani and S{\o}rensen(2008)]{ottaviani2008}
Ottaviani, M., S{\o}rensen, P.N., 2008.
The favorite-longshot bias: an overview of the main explanations.
In: Hausch, D.B., Ziemba, W.T. (Eds.), \emph{Handbook of Sports and Lottery
Markets.} Elsevier, Amsterdam, pp. 83--101.

\bibitem[Qin and Yang(2026)]{qin2026}
Qin, B., Yang, R., 2026.
\emph{Polymarket-v1 Database.} arXiv:2606.04217.

\bibitem[Simpson(1951)]{simpson1951}
Simpson, E.H., 1951.
The interpretation of interaction in contingency tables.
\emph{Journal of the Royal Statistical Society, Series B} 13 (2), 238--241.

\bibitem[Snowberg and Wolfers(2010)]{snowberg2010}
Snowberg, E., Wolfers, J., 2010.
Explaining the favorite--long shot bias: is it risk-love or misperceptions?
\emph{Journal of Political Economy} 118 (4), 723--746.

\bibitem[Sun and Abraham(2021)]{sun2021}
Sun, L., Abraham, S., 2021.
Estimating dynamic treatment effects in event studies with heterogeneous
treatment effects.
\emph{Journal of Econometrics} 225 (2), 175--199.

\bibitem[Thaler and Ziemba(1988)]{thaler1988}
Thaler, R.H., Ziemba, W.T., 1988.
Anomalies: parimutuel betting markets: racetracks and lotteries.
\emph{Journal of Economic Perspectives} 2 (2), 161--174.

\bibitem[Whelan(2023)]{whelan2023}
Whelan, K., 2023.
\emph{How Do Prediction Market Fees Affect Prices and Participants?}
CEPR Discussion Paper 17972; MPRA Paper 116926, University Library of Munich.

\bibitem[Wolfers and Zitzewitz(2004)]{wolfers2004}
Wolfers, J., Zitzewitz, E., 2004.
Prediction markets.
\emph{Journal of Economic Perspectives} 18 (2), 107--126.

\end{thebibliography}

\end{document}
